\section{Experimental Methodology}
\label{sec:method}

Simulations use ns-3.48~\cite{ns3} with AODV routing over IEEE 802.11.
Network parameters appear in Table~\ref{tab:sim} and detector parameters in
Table~\ref{tab:det}; both were fixed for the entire campaign and no per-scenario
tuning was performed.

\begin{table}[!t]
\caption{Network and simulation parameters.}
\label{tab:sim}
\centering
\small
\begin{tabular}{@{}llr@{}}
\toprule
Parameter & Symbol & Value \\
\midrule
Simulator & --- & ns-3.48 \\
Routing protocol & --- & AODV \\
Field & --- & $500 \times 500$\,m \\
Nodes & $N$ & 25 \\
Monitored beacons & $|\mathcal{B}|$ & 3 (1 attacker) \\
Transmit power & --- & 30\,dBm \\
Propagation loss model & --- & log-distance \\
Path-loss exponent & $\eta$ & 3.0 \\
Reference loss at 1\,m & $P_0$ & 46.68\,dB \\
Shadowing & $\sigma_{sh}$ & 0\,dB; 4 and 8\,dB in Sec.~\ref{sec:shadow} \\
RSSI measurement noise & $\sigma_{RSSI}$ & 4\,dB \\
Mobility model & --- & random waypoint \\
Node speed & --- & $U(1,5)$\,m/s \\
Pause time & --- & 2.0\,s \\
Traffic & --- & 10 CBR flows, 12\,pkt/s \\
Packet size & --- & 512\,B \\
Beacon interval & --- & 1.0\,s \\
PHY data rate & --- & DSSS 1\,Mbps \\
Simulation duration & $T$ & 400\,s \\
Attack start & --- & $t = 30$\,s \\
Per-axis position offset (A1) & $\delta$ & 300\,m \\
Measured mean displacement (A1) & --- & 267\,m \\
Drop probability (A2) & $p$ & 0.6 \\
Hybrid scaling (A3) & --- & $\delta/2$, $p/2$ \\
Runs executed / distinct & --- & 570 / 525 (Table~\ref{tab:campaign}) \\
\bottomrule
\end{tabular}
\end{table}

\begin{table}[!t]
\caption{Detector parameters. All values fixed across the campaign; none were
tuned per scenario.}
\label{tab:det}
\centering
\small
\begin{tabular}{@{}llr@{}}
\toprule
Parameter & Symbol & Value \\
\midrule
Observation window / step & --- & 10\,s / 5\,s \\
Branch weights & $\omega_B$, $\omega_L$ & 1.0, 1.0 \\
Confidence half-weight, localisation & $n^L_{1/2}$ & 5 \\
Confidence half-weight, behaviour & $n^B_{1/2}$ & 45 \\
Minimum peer set & $n_{min}$ & 8 \\
Minimum transit packets, behaviour & --- & 5 \\
Trust weights & $w_D$, $w_I$, $w_H$ & 0.5, 0.2, 0.3 \\
Trim fraction / recommender cap & --- / $\kappa$ & 0.2 / 0.3 \\
History decay, falling / rising & $\lambda_\downarrow$ / $\lambda_\uparrow$ & 0.5 / 0.9 \\
Observer agreement & $a_{min}$ & 0.5 \\
Malicious threshold & $\tau_M$ & 0.40 ($|Z| = 1.354$) \\
Normal threshold & $\tau_N$ & 0.70 ($|Z| = 0.845$) \\
Persistence / clearing persistence & $K$ / $K_c$ & 3 / 12 windows \\
Minimum observers & --- & 3 \\
Single-branch override & $\tau_X$ & 0.10 ($|Z| = 2.146$) \\
Bounded transform & $\phi$ & $\exp(-Z^2/2)$, Eq.~\eqref{eq:phi} \\
\bottomrule
\end{tabular}
\end{table}

Two properties of the radio configuration bound what the localisation branch
is asked to do. First, the primary campaign disables shadowing, so the channel
is deterministic log-distance path loss and the only stochasticity entering
the localisation branch is the 4\,dB measurement noise added at the receiver.
A real deployment would add multipath, obstruction and correlated shadowing,
all of which inflate the range residual for honest beacons and raise the
detection floor; Section~\ref{sec:shadow} repeats the campaign at
$\sigma_{sh} = 4$ and 8\,dB. Second, the observer inverts distance from
measured RSSI using the same path-loss parameters,
\begin{equation}
\label{eq:dhat}
\hat d_{ij} = 10^{(P_{tx} - P_0 - \mathrm{RSSI}_{ij})/(10\eta)},
\end{equation}
so $\hat d_{ij}$ in~\eqref{eq:resid} is derived entirely from a quantity the
observer measures. No true position enters the detector.

The full-intensity attack magnitudes are large. A position offset of
$\delta = 300$\,m in a $500 \times 500$\,m field is a very large lie, and a
drop probability of $p = 0.6$ is an aggressive grey hole. Both were chosen to
establish the mechanism rather than to represent a realistic adversary;
Section~\ref{sec:sweep} sweeps the hybrid attack down to a displacement of
3.75\,m and a drop probability of 0.75\%.

\begin{table}[!t]
\caption{Simulation campaign. Calibration seeds are clean runs held out from
testing. Of the 570 runs, 45 repeat a configuration and seed run elsewhere in
the campaign and reproduced it byte for byte; 525 are distinct.}
\label{tab:campaign}
\centering
\small
\begin{tabular}{@{}p{3.3cm}p{6.6cm}r@{}}
\toprule
Part & Design & Runs \\
\midrule
Original campaign & 4 scenarios $\times$ 5 seeds (detection) and 3 attacks
$\times$ 3 policies $\times$ 5 seeds with isolation enforced in the simulator & 65 \\
30-seed campaign & 4 scenarios $\times$ seeds 1--30, plus calibration seeds
101--105 & 125 \\
Intensity sweep, primary & A3, 11 levels $\times$ 10 seeds, drop $= 0.3f$ & 110 \\
Intensity sweep, second pairing & A3, 8 levels $\times$ 10 seeds, drop $= 0.15f$ & 80 \\
Shadowing & $\sigma_{sh} \in \{4, 8\}$\,dB, 4 scenarios $\times$ 20 seeds, plus
5 calibration seeds each & 170 \\
Clean enforced replay & policies D$_1$, E on clean seeds 1--10, schedules
from the replay enforced in the simulator & 20 \\
\midrule
Total & & 570 \\
\bottomrule
\end{tabular}
\end{table}

Table~\ref{tab:campaign} lists the campaign. Three groups of runs repeat
another part of it: the control arm of the enforced runs (15 runs) repeats the
original detection runs, those detection runs (20) are seeds 1--5 of the
30-seed campaign, and the primary sweep's full-intensity level (10) is seeds
1--10 of the 30-seed A3 scenario. Every repeated run reproduced its
counterpart's output files exactly, which we report as a determinism check;
no analysis uses the same simulation twice. Throughput recovery, mechanism
accounting, localisation error and the aggregation ablation come from the
original campaign, in which isolation is enforced by the simulator. Coverage,
detection latency, false positives and cumulative false isolation come from
the 30-seed campaign. For the latter, each isolation policy is applied to the
same per-window verdict stream offline: $\Phi$ depends only on when each beacon
is isolated and released, so one detection run per seed yields every policy's
$\Phi$ with the seeds paired across policies. The replay assumes that an
isolation does not alter the verdicts issued after it. How far that assumption
reaches differs by policy. D, D$_w$ and D$_1$ never release, so the
contribution of a beacon to $\Phi$ is fixed by its first isolation time, and
that time is computed exactly from verdicts issued before any isolation
occurred; feedback can shift only the isolation times of beacons removed
later in the same run. Policy E releases on a revised verdict, so its release
times, and hence its $\Phi$, depend on verdicts issued while another beacon may
be isolated. Removing a beacon from forwarding shifts load onto its neighbours
and can move their peer-relative scores in either direction. The enforced
runs of the original campaign are all attacked runs and do not test the
clean-network replay, so we tested it separately: for clean seeds 1--10 the
replayed D$_1$ and E schedules were enforced in the simulator, the enforced
traces were re-verified with the same calibration, and each policy's schedule
was rebuilt from the verdicts produced under isolation (20 further runs).
This is one step of a fixed-point check rather than a fully online detector,
and Section~\ref{sec:phi} reports the result. We compare six isolation
policies:

\begin{itemize}
\item \textbf{C, no isolation.} No node removed. Control.
\item \textbf{D, immediate isolation.} Isolate on the first window in which a
      beacon is Suspicious or Confirmed; never release.
\item \textbf{D$_w$, immediate with warm-up guard.} As D, but a beacon's first
      three usable windows cannot trigger isolation. The guard is counted per
      beacon from its own first usable window, never from the attack start.
\item \textbf{D$_1$, single window of malicious evidence.} Isolate on the
      first window in which confidence falls below $\tau_M$ with enough
      observers in agreement, without the persistence requirement; never
      release. This is the strongest honest immediate policy.
\item \textbf{E, confirmation-gated isolation.} Isolate only on confirmation;
      release when the verdict is subsequently revised.
\item \textbf{E$'$, confirmation-gated isolation with evidence-aware
      release.} As E, but a beacon confirmed while behavioural evidence was
      available is released only when that evidence returns
      (Section~\ref{sec:framework}). E$'$ was introduced after the enforced
      runs showed that E re-admits forwarding attackers
      (Section~\ref{sec:release}); because a beacon confirmed on behaviour
      stays isolated under E$'$, its closed-loop schedule follows exactly
      from the verdicts issued before isolation.
\end{itemize}

Policies D, D$_w$ and D$_1$ are constructions for comparison and are not models
of any published method; in particular they are not EMBN-MANET, whose response
operates on localisation coordination rather than packet forwarding.
Isolation is enforced by a blacklist queueing discipline installed on every
node, including in control runs. With an empty blacklist the discipline
behaves as a plain FIFO, so its queueing behaviour is present in the baseline
as well.

\subsection{Reproducibility}

The simulator is ns-3.48 built with GCC on Ubuntu under WSL2, with the analysis
pipeline in Python~3. Simulator sources, the analysis pipeline and the scripts
for every campaign are public at
\url{https://github.com/kraghu356/manet-beacon-trust}. Every run uses an explicit integer seed, so any run is
reproducible from its scenario, policy and seed. Detection thresholds were
fixed on the trust scale before any attack scenario was run, and the null
distributions of $A_B$ and $A_L$ were calibrated on clean-network runs that
are never used for testing: seeds 101--105 for the 30-seed and shadowing
campaigns, recalibrated per channel. No threshold was chosen by reference to a
result reported in this paper. Several simulator parameters are supplied at
invocation rather than compiled in, so the released command lines, not the
source defaults, define the experiment. Statistical intervals are percentile
bootstrap intervals over seeds (10{,}000 resamples), paired across policies
where policies share seeds.
